%% =====================================================================
%% pai-ban / assets/preamble.tex   —  版式实现（唯一权威规格：rules/design-spec.md）
%%
%% 用法：pandoc -H preamble.tex（由 scripts/build.py 复制到工作目录，
%%       替换 @...@ 占位符后使用）。pandoc 的默认模板已经把 xcolor /
%%       amsmath / unicode-math / microtype / longtable+booktabs+array /
%%       caption / calc / etoolbox / hyperref / xeCJK 全部加载完毕，
%%       本文件位于 header-includes（模板最末），只做"改"，不重复 load。
%%
%% 硬约束（实测 TeX Live 2026basic）：不得使用 titlesec / tocloft /
%%       tcolorbox / fvextra / framed / enumitem / adjustbox / mdframed /
%%       needspace（均不在 basic 里）。本文件用到的每个宏包见下，
%%       全部 kpsewhich 可查到。
%%
%% 占位符（build.py 替换；残留 @...@ 会让 build.py 直接报错）：
%%   @BOOKTITLE@   书名（偶数页页眉）
%%   @MAINFONT@    西文正文（Times New Roman）
%%   @MONOFONT@    等宽（Menlo，须覆盖 ░ █ ─ ├ └ │）
%%   @CJKFONT@     中文正文（Songti SC）
%%   @HEADFONT@    中文标题（PingFang SC）
%%   @SYMBOLFONT@  符号回退（Apple Symbols，⚠ U+26A0）
%%   @TOCNAME@     目录标题（目录 / Contents）
%%   @PARINDENT@   段首缩进（中文 2em / 西文 1.5em）
%%   @BOOKLANG@    语言开关：zh | en（影响引号与中西文间距策略）
%% =====================================================================

\usepackage{fancyhdr}      % 页眉页脚（§6）
\usepackage{ragged2e}      % \RaggedRight，避免 \raggedright 的坏断行
\usepackage{etoolbox}      % \AtBeginEnvironment / \patchcmd（§8 §10）
\usepackage{xstring}       % 字符串工具（未用，保留以备审计扩展）
\usepackage{setspace}      % 备用；正文行距由字号阶梯直接给定（§2.2）
\usepackage{graphicx}      % 图（§9）。pandoc 只在"文档里有图"时才自动加载
                           % graphicx，而本技能的图被 figure-caption.lua 转成
                           % raw LaTeX 后 pandoc 统计不到 → 必须自己加载
\usepackage[figuresright]{rotating}  % §9 第三级：整页横向页（sidewaysfigure）。
                           % 实测 figuresright = 图顶朝纸张左缘、题注落在右侧，
                           % 即读者顺时针转书即可阅读；全书方向一致（不随奇偶
                           % 页翻转，否则同一本书要左转一次右转一次）。
\usepackage{floatpag}      % \thisfloatpagestyle：旋转页必须"无页眉、页码正常"
                           % （§6/§9）。\thispagestyle 对浮动页无效，实测页眉
                           % 仍会印出书名/章名。
\usepackage{listings}      % 代码块（§7）
\usepackage{booktabs}      % 三线表（§8）
\usepackage{longtable}     % 跨页长表（§8）
\usepackage{tabularx}      % 后备表格
\usepackage{array}         % 列格式 >{}
\usepackage{calc}          % p{} 列宽算式
\usepackage{colortbl}      % 表头底色（§8，默认不用）
\usepackage{caption}       % 题注格式（§8 §9）
\usepackage{bookmark}      % PDF 书签层级（§5）

%% ---------------------------------------------------------------------
%% spec §3 色彩：全书只有一种强调色 + 灰阶（颜色种类 ≤ 6）
%%   灰度换算（用于 §3 的黑白打印验收）：
%%     ink    #1A1A1A → 10%   accent #0B3D91 → 24%
%%     两者差 14%，但强调色只用于章号/链接/关键字等**非正文**元素，
%%     正文与强调色的亮度差在 25% 以上的是 muted→ink 与 rule→ink 两对。
%% ---------------------------------------------------------------------
\definecolor{ink}{HTML}{1A1A1A}        % 正文
\definecolor{accent}{HTML}{0B3D91}     % 章号 / 链接 / 代码关键字
\definecolor{accent30}{HTML}{B6C5DE}   % accent 30% + 白 70%（章号大字）
\definecolor{rule}{HTML}{D0D4DA}       % 分隔线 / 页眉线
\definecolor{muted}{HTML}{6E7781}      % 页码 / 代码行号 / 注释 / 代码标签
\definecolor{codebg}{HTML}{F7F8FA}     % 代码块底色
\definecolor{strgreen}{HTML}{1B7F3B}   % 代码字符串（§7，不计入"书籍色彩"）
\definecolor{numrust}{HTML}{8A3B12}    % 代码数字（§7）
%% 注意：strgreen / numrust 是 §7 明文规定的代码语法色，属于印刷惯例而非
%% 版面色。§3 的"颜色 ≤ 6"验收只统计版面语义色（ink/accent/rule/muted/
%% codebg/tablehead），verify.py 按这 6 个统计。

%% ---------------------------------------------------------------------
%% spec §2.1 字体族
%%   * 正文西文由 pandoc 的 -V mainfont 设好（Times New Roman）；
%%   * 中文正文由 -V CJKmainfont 设好（Songti SC）；
%%   * 这里补：中文标题族（PingFang SC Semibold）、等宽中文字形回退、
%%     符号回退字体。
%% ---------------------------------------------------------------------
\newfontfamily\symbolfallback{@SYMBOLFONT@}
%% ⚠ (U+26A0)：Times / Songti 均无字形 → 走 Apple Symbols（§12 第 1 步的
%% 缺字形替换把它变成 `\warn`{=latex}，见 scripts/build.py）
\newcommand{\warn}{{\symbolfallback\char"26A0\relax}}

%% 中文标题族：PingFang SC 的 Bold 面在 macOS 上叫 Semibold，
%% 不显式指定 BoldFont 会触发 fontspec 的 "Font shape undefined" 警告
%% （不会报错，但标题会掉回 Regular，字重不够）。
\setCJKfamilyfont{heiti}[BoldFont={PingFang SC Semibold},Scale=1]{@HEADFONT@}
\providecommand{\heiti}{\CJKfamily{heiti}}

\makeatletter
\@ifpackageloaded{xeCJK}{%
  %% 代码块里的中文注释（§7）：\ttfamily 下的 CJK 走 CJKmonofont
  \setCJKmonofont{@CJKFONT@}%
  \setCJKsansfont{@CJKFONT@}%
  %% 圈码 ①-⑮（U+2460-246F）落在 xeCJK 默认 CJK 范围之外，
  %% 不声明就会掉到 Times → Missing character
  \xeCJKDeclareCharClass{CJK}{"2460 -> "246F}%
  %% 制表符 / 方块元素（░ █ ─ ├ └ │）：Menlo 有字形，但它们不是 CJK，
  %% 落在主字体里就会缺字。把几何图形区（U+2500-25FF 制表符）声明为
  %% "Default" 以外的类没有意义——正确做法是用 \newfontfamily 包一层，
  %% 见下方 \paibangeo。
}{}
\makeatother

%% 制表符 / 方块元素回退（正文里出现时用；代码块内由 Menlo 直接覆盖）
\newfontfamily\geofont{@MONOFONT@}
\newcommand{\paibangeo}[1]{{\geofont #1}}

%% ---------------------------------------------------------------------
%% spec §2.2 字号阶梯（严格 1.25 模数）
%%   章标题 22/26 Bold      节 14/17 Bold      小节 11.5/14 Bold
%%   三级 11/13.5 BoldItal  正文 11/15.2       表格 9/11.5
%%   图题 9/11.5            代码 8.5/10.5      引文 9.5/13
%%   脚注 8/10              页眉页脚 8
%%   实现方式：直接重定义字号宏，让每一级都落在阶梯上（不靠 linespread
%%   的连乘——11pt 类的 \baselineskip 是 13.6pt，linespread 1.12 得
%%   15.23pt 是"差不多"，而 §0 原则 1 要求精确）。
%% ---------------------------------------------------------------------
\makeatletter
\renewcommand\normalsize{\@setfontsize\normalsize{11}{15.2}}   % 正文
\renewcommand\small{\@setfontsize\small{9}{11.5}}              % 表格 / 题注
\renewcommand\footnotesize{\@setfontsize\footnotesize{8}{10}}  % 脚注 / 页眉
\renewcommand\scriptsize{\@setfontsize\scriptsize{7}{8.5}}     % 代码行号
\renewcommand\large{\@setfontsize\large{11.5}{14}}             % 小节
\renewcommand\Large{\@setfontsize\Large{14}{17}}               % 节
\renewcommand\LARGE{\@setfontsize\LARGE{22}{26}}               % 章
\renewcommand\Huge{\@setfontsize\Huge{48}{52}}                 % 章号（§4）
\makeatother

%% ---------------------------------------------------------------------
%% spec §0 原则 1 / §2.2 正文：段首缩进、段间 0、基线与行距
%%   pandoc 默认可能加载 parskip.sty（会把 \parindent 清零），这里必须在
%%   它之后重新压回来。
%% ---------------------------------------------------------------------
\setlength{\parindent}{@PARINDENT@}
\setlength{\parskip}{0pt}
\setlength{\topskip}{11pt}           % 每页首行基线 = 版心顶 + 11pt
\clubpenalty=10000                   % 不留孤行（§0 原则 6）
\widowpenalty=10000
\displaywidowpenalty=10000
\brokenpenalty=10000
\setlength{\emergencystretch}{6em}   % 长 URL 的最后一道闸（§7/§11；4em 时
                                     % modelcontextprotocol.io 那一行仍溢出 15.8pt，
                                     % 6em 是实测能压住的最小值）
\raggedbottom                        % 双面书：宁可底部留白，不撑散行距

%% ---------------------------------------------------------------------
%% spec §1 画布与版心
%%   A5 148×210，双面；内 18 / 外 15 / 上 16（含页眉区）/ 下 19；
%%   headsep 5mm，footskip 7mm。
%%   geometry 由 build.py 用 -V geometry=... 传入，这里只做校验性断言：
%%   版心必须是 115×175mm（±0.5mm），否则报错退出（宁可停，不出错版）。
%% ---------------------------------------------------------------------
\makeatletter
\newlength{\paiban@textw}\setlength{\paiban@textw}{\textwidth}
\newlength{\paiban@texth}\setlength{\paiban@texth}{\textheight}
\AtBeginDocument{%
  \ifdim\paiban@textw<324pt \errmessage{pai-ban: text width < 114mm, check geometry}\fi
  \ifdim\paiban@textw>329pt \errmessage{pai-ban: text width > 116mm, check geometry}\fi
  \ifdim\paiban@texth<494pt \errmessage{pai-ban: text height < 174mm, check geometry}\fi
  \ifdim\paiban@texth>498pt \errmessage{pai-ban: text height > 176mm, check geometry}\fi
}
\makeatother

%% ---------------------------------------------------------------------
%% spec §4 章首页（签名元素）
%%   起于奇数页（\cleardoublepage）→ 24mm 空白 → 章号 48pt（accent 30%）
%%   → 6pt → 0.8pt accent 短线 24mm → 8pt → 章标题 22/26 Bold
%%   → 18mm → 正文。章首页 plain（无页眉，页码居中页脚）。
%%   编号与标题由 chapter-split.lua 拆开传入：\paibanchapter{2}{标题}
%%   无编号章节：\paibanchapter*{前言}
%% ---------------------------------------------------------------------

%% 空白页也要有页码（§6：空白页 plain，不排页眉），否则读者会以为漏印
\makeatletter
\renewcommand{\cleardoublepage}{%
  \clearpage
  \if@twoside
    \ifodd\c@page\else
      \hbox{}\thispagestyle{plain}\newpage
      \if@twocolumn\hbox{}\newpage\fi
    \fi
  \fi}

%% pandoc 的 book 模板会无条件在正文前插入 \mainmatter（阿拉伯页码从 1 起）。
%% 我们需要「前言/目录 = 罗马数字，第一章 = 阿拉伯数字 1」，所以把模板那次
%% \mainmatter 变成空操作，改由 chapter-split.lua 在**第一个编号章**前
%% 发出 \paibanmainmatter。
\let\paiban@realmainmatter\mainmatter
\renewcommand{\mainmatter}{\relax}
\newcommand{\paibanmainmatter}{\paiban@realmainmatter}

\newcommand{\paiban@chapmark}{}
\newcommand{\paibanbooktitle}{@BOOKTITLE@}

%% 章首页版式：#1 = 章号（可空），#2 = 标题（已是 LaTeX 片段）
\newcommand{\paiban@opener}[2]{%
  \vspace*{24mm}%
  \if\relax\detokenize{#1}\relax\else
    {\fontsize{48}{52}\selectfont\color{accent30}\bfseries #1\par}%
    \vspace{6pt}%
    {\color{accent}\rule{24mm}{0.8pt}\par}%
    \vspace{8pt}%
  \fi
  {\RaggedRight\fontsize{22}{26}\selectfont\heiti\bfseries\color{ink}%
   #2\par}%
  \vspace{18mm}%
}

%% 编号章
\providecommand{\paibannumlabel}[1]{@CHAPNUMBEFORE@#1@CHAPNUMAFTER@}%  目录条目的章号样式（zh: 第N章 / en: Chapter N:）
\newcommand{\paiban@chapnum}[2]{%
  \cleardoublepage
  \thispagestyle{plain}%          §6 章首页 plain
  \global\@topnum\z@
  \phantomsection                 % hyperref 锚点（否则目录链接跳到上一章）
  \addcontentsline{toc}{chapter}{\paibannumlabel{#1}\hspace*{0.75em}#2}%
  \markboth{#2}{#2}%              §6 奇数页页眉 = 当前章标题
  \gdef\paiban@chapmark{#2}%
  \paiban@opener{#1}{#2}%
}

%% 无编号章（前言 / 参考文献 / 附录）：无章号、无短线，标题 22pt
\newcommand{\paiban@chapstar}[1]{%
  \cleardoublepage
  \thispagestyle{plain}%
  \global\@topnum\z@
  \phantomsection
  \addcontentsline{toc}{chapter}{#1}%
  \markboth{#1}{#1}%
  \gdef\paiban@chapmark{#1}%
  \paiban@opener{}{#1}%
}

\newcommand{\paibanchapter}{\@ifstar{\paiban@chapstar}{\paiban@chapnum}}

%% ---------------------------------------------------------------------
%% spec §2.2 节 / 小节 / 三级标题
%%   用 \@startsection 重定义（不用 titlesec）。正文里节号（"2.1"）是源文
%%   自带的文本，所以这里不排自动编号（secnumdepth 已被 pandoc 设为
%%   负无穷），目录条目手工 \addcontentsline。
%%   beforeskip 取正值 → 标题后首段仍缩进（中文书惯例）。
%% ---------------------------------------------------------------------
\newcommand{\paiban@sectionstyle}{%
  \fontsize{14}{17}\selectfont\heiti\bfseries\color{ink}}
\newcommand{\paiban@subsectionstyle}{%
  \fontsize{11.5}{14}\selectfont\heiti\bfseries\color{ink}}
\newcommand{\paiban@subsubsectionstyle}{%
  \fontsize{11}{13.5}\selectfont\heiti\bfseries\itshape\color{ink}}

%% 两个必须注意的地方（都踩过）：
%%   1. 标题必须用 `{#1}` 传，不能写成 `#1`。\@startsection 的第 7 个参数
%%      由 \@dblarg 读取，它只看**下一个 token**：`#1` 展开后不带花括号，
%%      于是只吃到 "2.1" 里的 "2"，标题后半截掉进正文（实测 TOC 里出现
%%      孤立条目「2」、标题排成 ".1 用Ollama…"）。
%%   2. **不要**自己 \addcontentsline。LaTeX 内核的 \@sect 已经按
%%      \c@tocdepth 写好目录条目；再写一次就是两条（一条正确、一条截断）。
%%      \phantomsection 仍要保留：它给内核那条目录条目提供 hyperref 锚点。
\newcommand{\paibansection}[1]{%
  \phantomsection
  \@startsection{section}{1}{\z@}{14pt}{6pt}{\paiban@sectionstyle}{#1}}
\newcommand{\paibansubsection}[1]{%
  \phantomsection
  \@startsection{subsection}{2}{\z@}{11pt}{4pt}{\paiban@subsectionstyle}{#1}}
\newcommand{\paibansubsubsection}[1]{%
  \phantomsection
  \@startsection{subsubsection}{3}{\z@}{9pt}{3pt}{\paiban@subsubsectionstyle}{#1}}

%% 章标题在页眉里不能带 LaTeX 命令；\markboth 收到的是干净的文本片段
%% （chapter-split.lua 用 pandoc.write 渲染，标题里若有 \texttt 会原样出现，
%% 这在 8pt 页眉里是可接受的）。

%% ---------------------------------------------------------------------
%% spec §5 目录
%%   * 只收 chapter / section（pandoc 的 toc-depth=2）
%%   * 章条目 11pt Bold，段前 6pt，引导点，页码右对齐
%%   * 节条目 9.5pt，左缩进 1.4em，引导点，页码右对齐
%%   * 目录页无页眉，罗马数字（\frontmatter 由 pandoc 模板发出）
%%   不用 tocloft：直接重写 \l@chapter / \l@section。
%% ---------------------------------------------------------------------
\renewcommand*{\l@chapter}[2]{%
  \ifnum \c@tocdepth >\m@ne
    \addpenalty{-\@highpenalty}%
    \vskip 6pt \@plus 2pt \@minus 2pt
    \begingroup
      \leftskip\z@ \rightskip\@pnumwidth \parfillskip-\@pnumwidth
      \noindent\fontsize{11}{15.2}\selectfont\bfseries\color{ink}%
      #1\nobreak\leaders\hbox{$\m@th\mkern\@dotsep mu\hbox{.}\mkern\@dotsep mu$}%
        \hfill\nobreak
      \hb@xt@\@pnumwidth{\hss\fontsize{9.5}{11.5}\selectfont\color{muted}#2}%
      \par
    \endgroup
    \penalty\@highpenalty
  \fi}
\renewcommand*{\l@section}[2]{%
  \ifnum \c@tocdepth >\z@
    \vskip 2pt
    \begingroup
      \leftskip1.4em \rightskip\@pnumwidth \parfillskip-\@pnumwidth
      \noindent\fontsize{9.5}{12}\selectfont\color{ink}%
      #1\nobreak\leaders\hbox{$\m@th\mkern\@dotsep mu\hbox{.}\mkern\@dotsep mu$}%
        \hfill\nobreak
      \hb@xt@\@pnumwidth{\hss\color{muted}#2}%
      \par
    \endgroup
  \fi}

%% 目录页：与章首页同一套骨架（24mm 顶空 + 22pt 标题 + 10mm），
%% 但不排页眉、页码居中（§5 §6）
\renewcommand{\tableofcontents}{%
  \cleardoublepage
  \thispagestyle{plain}%
  \phantomsection
  \pdfbookmark[0]{\contentsname}{toc}%
  \vspace*{24mm}%
  {\RaggedRight\fontsize{22}{26}\selectfont\heiti\bfseries\color{ink}\contentsname\par}%
  \vspace{10mm}%
  {\pagestyle{plain}\@starttoc{toc}}%   目录独占 plain（后续页也不排页眉，§5）
}
\makeatother

%% ---------------------------------------------------------------------
%% spec §6 页眉与页码
%%   双面：偶数页（verso）页眉 = 书名（居外 = 左），
%%         奇数页（recto）页眉 = 当前章标题（居外 = 右）
%%   8pt muted，字距 +0.05em（\textls[50] = 50/1000 em），
%%   页眉下 0.4pt rule 细线，长度 = 版心宽；
%%   页码全书统一放页脚居中（§6 只能选一种）。
%% ---------------------------------------------------------------------
\pagestyle{fancy}
\fancyhf{}
\fancyhead[LE]{\fontsize{8}{10}\selectfont\color{muted}\textls[50]{\paibanbooktitle}}
\fancyhead[RO]{\fontsize{8}{10}\selectfont\color{muted}\textls[50]{\leftmark}}
\fancyfoot[C]{\fontsize{8}{10}\selectfont\color{muted}\thepage}
\renewcommand{\headrulewidth}{0.4pt}
\renewcommand{\footrulewidth}{0pt}
%% fancyhdr 的 \headrule 用当前颜色画线 → 单独重画成 rule 色
\renewcommand{\headrule}{\color{rule}%
  \hrule width\headwidth height\headrulewidth\relax
  \vskip-\headrulewidth}
%% plain：章首页 / 目录页 / 空白页（无页眉、无页眉线、页码居中）
\fancypagestyle{plain}{%
  \fancyhf{}%
  \fancyfoot[C]{\fontsize{8}{10}\selectfont\color{muted}\thepage}%
  \renewcommand{\headrulewidth}{0pt}%
  \renewcommand{\headrule}{}%
}
%% 浮动体页（图/表独占页）沿用 fancy，但也要页眉线 —— 默认继承，无需处理

%% ---------------------------------------------------------------------
%% spec §7 代码块（listings；禁止 minted / tcolorbox / fvextra）
%%   浅灰底 codebg + 左侧 2pt accent 竖线 + 内边距 4pt + 无外框；
%%   Menlo 8.5/10.5；行号仅源文有 linenums 时出现（左侧 muted 7pt）；
%%   breaklines 折行，续行前缀 ↪（muted）；
%%   语法四色：关键字 accent / 字符串墨绿 / 注释 muted 斜体 / 数字赭石。
%% ---------------------------------------------------------------------
%% 把 pandoc 的 language=<name> 映射到 listings 认识的方言。
%% listings 没有 bash 方言（有 sh），没有 console/plaintext；不 alias 的话
%% 会 "Couldn't load requested language" 而把整块代码排成纯文本。
\lstalias{bash}{sh}
\lstalias{console}{sh}
\lstalias{shell}{sh}
\lstalias{zsh}{sh}
\lstalias{text}{}
\lstalias{txt}{}
\lstalias{plaintext}{}

\lstdefinestyle{paiban}{%
  backgroundcolor=\color{codebg},%      §7 浅灰底
  frame=l, %                             只画左边一条
  framerule=2pt,%                        2pt 强调色竖线
  rulecolor=\color{accent},%
  framesep=4pt,%                         §7 内边距 4pt
  xleftmargin=20pt,%                     整体右移：给行号留出内侧空档
  xrightmargin=2pt,%
  framexleftmargin=12pt,%                竖线左移 12pt，把行号包进灰底内
  %% 为什么是 20/12：行号右端在 xleftmargin−numbersep=14pt，三位数在
  %% Menlo 7pt 下宽 10.5pt → 左端距版心 3.5pt（实测 p.8「130」= 46.01 /
  %% verso 42.52），仍在版心内（§1 左边界）。
  %% 早期 xleftmargin=10/framexleftmargin=0 时行号挂在竖线**左边**、
  %% 两位数就伸进外白边 2.9pt，三位数必然越界（实测 p.28/p.42/p.48）。
  basicstyle=\fontsize{8.5}{10.5}\selectfont\ttfamily\color{ink},%   §2.2 代码 8.5/10.5
  keywordstyle=\color{accent}\bfseries,%
  stringstyle=\color{strgreen},%
  commentstyle=\color{muted}\itshape,%
  numberstyle=\fontsize{7}{8.5}\selectfont\color{muted},%
  numbers=none,%                         默认不显示；pandoc 传 numbers=left 时才显示
  numbersep=6pt,%
  breaklines=true,%                      §7 折行
  breakatwhitespace=false,%              中文注释行没有空格，必须允许任意位置断
  breakindent=0pt,%
  postbreak=\mbox{\textcolor{muted}{\scriptsize$\hookrightarrow$}\space},%
  columns=fullflexible,%
  keepspaces=true,%                      保留缩进（代码的可读性全靠它）
  showstringspaces=false,%
  tabsize=4,%
  extendedchars=true,%
  upquote=true,%
  aboveskip=8pt,%                        §2.2 块上下 8pt
  belowskip=8pt,%
  captionpos=t,%
  escapeinside={},%
}
\lstset{style=paiban}   % 全局生效：所有 lstlisting 都走 §7 的版式

%% listings 在 XeLaTeX 下处理 CJK 时会丢失"字符间可断行"的信息，
%% 这里显式把 CJK 标点与汉字之间的断点交回 xeCJK（保持默认即可），
%% 只需确保 extendedchars 打开、basicstyle 里已有 CJK 回退（CJKmonofont）。

%% --listings 会把行内代码变成 \passthrough{\lstinline!…!}。全局
%% breaklines=true 且 postbreak 里含 \mbox 时，一旦 TeX 正好在行内代码
%% 边界换行，\lst@@discretionary 里的 \mbox 会触发
%% "Improper discretionary list" 并中止编译（实测 chapter-02）。
%% 行内代码需要的只是"可断行"（否则 search_product(query, category)
%% 这种长串在 115mm 栏宽里必然 Overfull），不需要 ↪ 前缀。
%% 因此在 \passthrough 的局部组里保留 breaklines、清空 postbreak。
\renewcommand{\passthrough}[1]{{\lstset{breaklines=true,postbreak={}}#1}}

%% 代码块上方的文件标签（§7：title 属性 → 8.5pt muted 单行标签，不加框）
\newcommand{\paibancodecaption}[1]{%
  \par\addvspace{6pt}\noindent
  {\fontsize{8.5}{10.5}\selectfont\color{muted}\ttfamily #1\par}%
  \nobreak\vspace{-2pt}\nobreak}

%% ---------------------------------------------------------------------
%% spec §8 表格：booktabs 三线表，顶/底 0.8pt，栏目线 0.4pt，禁止竖线；
%%   9pt/11.5pt，tabcolsep 4pt；题注在表上方，"表 2.1" 加粗 + 空格 + 标题。
%% ---------------------------------------------------------------------
\setlength{\heavyrulewidth}{0.8pt}
\setlength{\lightrulewidth}{0.4pt}
\setlength{\cmidrulewidth}{0.4pt}
\setlength{\arrayrulewidth}{0.4pt}
\renewcommand{\arraystretch}{1.0}

\AtBeginEnvironment{longtable}{%
  \fontsize{9}{11.5}\selectfont
  \setlength{\tabcolsep}{4pt}%
  \renewcommand{\arraystretch}{1.0}%
}
\AtBeginEnvironment{tabular}{%
  \fontsize{9}{11.5}\selectfont
  \setlength{\tabcolsep}{4pt}%
}

%% 题注：pandoc 默认 "表 2.1: 标题"（冒号 + 正体标签）；
%% §8 要求 "表 2.1  标题"（加粗标签 + 空格，无冒号）。
\captionsetup[table]{%
  name=表,%
  labelsep=quad,%            加粗标签与标题之间一个 quad 的空格
  labelfont={bf,color=ink},%
  textfont={color=ink},%
  font={footnotesize},%      9/11.5
  position=top,%
  skip=4pt,%
  justification=raggedright,%
  singlelinecheck=false,%
  belowskip=6pt,%
}
\captionsetup[figure]{%
  name=图,%
  labelsep=quad,%
  labelfont={bf,color=ink},%
  textfont={color=ink},%
  font={footnotesize},%
  position=bottom,%
  skip=4pt,%
  justification=raggedright,%
  singlelinecheck=false,%
}

%% §8 表题注：table-caption.lua 把源文里「表14.1：…」的斜体段搬到表格
%% 上方，交给这个宏排版：#1 = 「表 14.1」（加粗），#2 = 标题。
%% 9pt / 11.5，与表间距 4pt；\nopagebreak 保证题注不会单独留在页尾。
\newcommand{\paibantablecaption}[2]{%
  \par\addvspace{6pt}%
  \noindent{\fontsize{9}{11.5}\selectfont
    {\bfseries\color{ink}#1\quad}#2\par}%
  \nopagebreak[4]\vspace{4pt}\nobreak}

%% 表头不单独留在页尾（§8）：longtable 的 \endhead 之后插一个惩罚，
%% 让"表头 + 至少一行"成为不可分块。
\AtBeginEnvironment{longtable}{\nopagebreak[3]}

%% ---------------------------------------------------------------------
%% spec §9 图：三级放置策略（策略与量测都在 scripts/figure_fit.py）
%%
%%   1) inline   裁边（trim+clip）后放进版心宽 115mm
%%   2) wide     裁边后仍 < 4.5pt → 加宽到 125mm（两侧各侵入白边 5mm，
%%               距纸边仍 ≥ 10mm；见 §9 与 verify.py 的门禁 2/3 例外）
%%   3) rotated  仍 < 4.5pt → 整页横向页（sidewaysfigure，可用宽 = 版心高 175mm）
%%
%% 谁决定用哪一级：build.py 调 figure_fit.py 量图（画布/墨迹 bbox/图内字号），
%% 逐图算出「成书后实际字号 = 图内字号 × 缩放比」，写进工作目录的
%% figure-fit.tex：
%%     \paibanfigdef{图路径}{inline|wide|rotated}{裁边 l b r t}{宽}{高}
%% 这里只负责查表执行。查不到（没有图清单）就退回第一版行为：不裁边、
%% 限宽限高——宁可少裁边，也不能因为清单丢了而不出书。
%%
%% 题注在图下方，9pt；图与题注必须同页：
%%   * inline/wide：不可分页的 minipage（图 + 题注 + 来源行）
%%   * rotated：sidewaysfigure 是浮动体，图与题注在同一个旋转盒里，
%%     题注跟着图一起转（这是"题注与图分离"的唯一可靠解法）
%% ---------------------------------------------------------------------
\newsavebox{\paibanfigbox}   % 也在 §10 里给 quote 量高度用
\newsavebox{\paibanpibox}
\newlength{\paibanpiw}
\newlength{\paibanpih}
\newlength{\paibanpiwmax}
\newlength{\paibanpihmax}

%% 放一张图：#1 = 图片路径。宽度上限 = 版心宽，高度上限 = 0.72 版心高，
%% 且**永不放大**超过原始尺寸（graphicx 同时给 width/height + keepaspectratio
%% 时按"装进这个盒子"缩放，所以 w_max = min(自然宽, 版心宽) 就够了）。
\newcommand{\paibaninclude}[1]{%
  \sbox{\paibanpibox}{\includegraphics{#1}}%
  \setlength{\paibanpiw}{\wd\paibanpibox}%
  \setlength{\paibanpih}{\dimexpr\ht\paibanpibox+\dp\paibanpibox\relax}%
  \setlength{\paibanpiwmax}{\linewidth}%
  \ifdim\paibanpiw<\paibanpiwmax \setlength{\paibanpiwmax}{\paibanpiw}\fi%
  \setlength{\paibanpihmax}{\dimexpr0.72\textheight\relax}%
  \includegraphics[width=\paibanpiwmax,height=\paibanpihmax,keepaspectratio]{#1}%
}

%% 单张图 + 题注的容器：#1 = 图片 LaTeX，#2 = 题注 LaTeX（已含"图 2.1"）
%% 用不可分页的 minipage 保证「图与题注同页」（§9），并让整块在换页时
%% 整体下移，绝不出现"图在页尾、题注在下一页"。
\newcommand{\paibanfigure}[2]{%
  \par\addvspace{10pt}%
  \noindent\begin{minipage}{\linewidth}%
    \centering
    #1\par
    \vspace{4pt}%                       §9 图与题注间距 4pt
    {\RaggedRight\fontsize{9}{11.5}\selectfont\color{ink}#2\par}%
  \end{minipage}%
  \par\addvspace{10pt}%
}

%% --- 放置表（build.py 生成）------------------------------------------
%% \paibanfigdef{路径}{方式}{裁边}{宽}{高} → 存成 \pbfig@<路径>（一个
%% 单参宏，参数是题注）。用路径做宏名：同一张图在书里出现两次时查同一行，
%% 不会因为题注不同而各放一次。
\newcommand{\paibanfigdef}[5]{%
  \expandafter\def\csname pbfig@#1\endcsname##1{%
    \csname paibanfig#2\endcsname{#1}{#3}{#4}{#5}{##1}}}
\InputIfFileExists{figure-fit.tex}{}{}

%% inline：图宽已由清单给定（≤ 版心宽），高度上限同时给，防长图撑爆页面
\newcommand{\paibanfiginline}[5]{% 路径 裁边 宽 高 题注
  \par\addvspace{10pt}%
  \noindent\begin{minipage}{\linewidth}%
    \centering
    \includegraphics[trim={#2},clip=true,width=#3,height=#4,keepaspectratio]{#1}\par
    \vspace{4pt}%
    {\RaggedRight\fontsize{9}{11.5}\selectfont\color{ink}#5\par}%
  \end{minipage}%
  \par\addvspace{10pt}%
}

%% wide：125mm 比版心宽 10mm，两侧各侵入 5mm。
%% 用 \makebox[0pt][c] 让"超出版心的部分"对称溢出，且不产生 Overfull
%% \hbox 警告（0pt 盒里的 \hss 吸收，§7/§8 的 Overfull 门禁不会误报）。
\newcommand{\paibanfigwide}[5]{%
  \par\addvspace{10pt}%
  \noindent\begin{minipage}{\linewidth}%
    \centering
    \makebox[0pt][c]{\includegraphics[trim={#2},clip=true,width=#3,height=#4,keepaspectratio]{#1}}\par
    \vspace{4pt}%
    {\RaggedRight\fontsize{9}{11.5}\selectfont\color{ink}#5\par}%
  \end{minipage}%
  \par\addvspace{10pt}%
}

%% rotated：整页横向页。
%%   * 图 + 题注在同一个旋转盒内 → 题注随图转，不可能分离；
%%   * \thisfloatpagestyle{plain} → 该页无页眉、页脚仍有页码（§6/§9）；
%%   * 紧跟 \clearpage → 浮动页立刻出，不会攒到章末（[p] 单独使用时
%%     LaTeX 只会在 \clearpage 时吐浮动页，攒到章末就是"图全在章尾"）。
\newcommand{\paibanfigrotated}[5]{%
  \par
  \begin{sidewaysfigure}[p]%
    \thisfloatpagestyle{plain}%
    \centering
    \includegraphics[trim={#2},clip=true,width=#3,height=#4,keepaspectratio]{#1}\par
    \vspace{4pt}%
    {\RaggedRight\fontsize{9}{11.5}\selectfont\color{ink}#5\par}%
  \end{sidewaysfigure}%
  \clearpage}

%% 过滤器调用入口：\paibanplaced{路径}{题注}
\newcommand{\paibanplaced}[2]{%
  \ifcsname pbfig@#1\endcsname
    \csname pbfig@#1\endcsname{#2}%
  \else
    \paibanfigure{\paibaninclude{#1}}{#2}%
  \fi}

%% 无题注的图（源文只放了图、没写斜体题注段）
\newcommand{\paibanplacedonly}[1]{%
  \ifcsname pbfig@#1\endcsname
    \begin{center}\csname pbfig@#1\endcsname{}\end{center}%
  \else
    \begin{center}\paibaninclude{#1}\end{center}%
  \fi}


%% 图题里"图 2.1"四个字加粗（§9）
\newcommand{\paibanfiglabel}[1]{{\bfseries #1}\hspace{0.6em}}

%% 图片来源行（> 数据来源：…）—— 8.5pt muted 斜体，紧跟题注
\newcommand{\paibansource}[1]{%
  \par\nobreak\vspace{2pt}%
  {\RaggedRight\fontsize{8.5}{10.5}\selectfont\color{muted}\itshape #1\par}}

%% pandoc 的 \pandocbounded：原先只做"限宽限高"，不居中。
%% 我们的 \paibanfigure 已把它放进 \centering 的 minipage，无需再包。
\makeatletter
\@ifundefined{pandocbounded}{}{%
  \let\paiban@pandocbounded\pandocbounded
  \renewcommand*\pandocbounded[1]{%
    {\centering\paiban@pandocbounded{#1}\par}}%
}
\makeatother

%% ---------------------------------------------------------------------
%% spec §10 引文 / 提示 / 列表
%%   引文块：左侧 2pt rule 竖线，左右缩进 1em，9.5/13，ink 80%
%%   提示块（> **注意**）：竖线用 accent
%%   实现：pandoc 的 blockquote → quote 环境；用 minipage 包一层拿"整块
%%   高度"来画竖线（不可分页——本技能面向的书稿里引文块都是 1–4 行的
%%   来源说明，宁可整块移到下一页，也不让竖线断成两截）。
%% ---------------------------------------------------------------------
%% 竖线：用 \vrule 撑满 minipage 的高度不可行（高度未知），
%% 改用 lrbox 量高度再画 \rule。**几何必须是"内容缩进 > 0"**：
%%   规格 §10 要求左右各缩进 1em，且 §1 要求内容不得侵入订口。
%%   早期实现用 \llap 把竖线挂在版心左缘之外（内容从版心 0 缩进开始，
%%   竖线落到订口里 9.8pt），被 §1 左边界门禁抓出来（12 页）。
%%   现在：整块右移 1em → 竖线落在 1em 处 → 0.8em 空隙 → 内容盒，
%%   内容盒宽 = 版心 − 1em(左) − (2pt + 0.8em) − 1em(右)，右缘仍留 1em。
%% \rule 只向上长，用 \raisebox[-\dp][0pt][0pt] 下移 dp，竖线才盖住
%% 最后一行降部；[0pt][0pt] 保证不影响行距。
\renewenvironment{quote}{%
  \par\addvspace{8pt}%
  \begin{lrbox}{\paibanfigbox}%
    \begin{minipage}{\dimexpr\linewidth-2.8em-2pt\relax}%
      \fontsize{9.5}{13}\selectfont\color{ink}%
      \setlength{\parindent}{0pt}%
      \RaggedRight
}{%
    \end{minipage}%
  \end{lrbox}%
  \noindent\hspace*{1em}%
  \raisebox{-\dp\paibanfigbox}[0pt][0pt]{%
    \textcolor{rule}{\rule{2pt}{\dimexpr\ht\paibanfigbox+\dp\paibanfigbox\relax}}}%
  \hspace{0.8em}%
  \usebox{\paibanfigbox}%
  \par\addvspace{8pt}%
}

%% 列表（§10）：项目间距 2pt，整体上下 4pt，缩进 1.5em；嵌套 1.2em
%% 注意（已知且刻意）：标签是右对齐到"体缩进 − \labelsep"，两位数
%% （参考文献的 19.）比一位数宽 5.5pt，会外挂约 2.8pt 进白边。
%% §10 规定 1.5em，§1 允许 3pt 容差，两者正好相容；三位数（100.）
%% 会外挂 8pt 并让 verify 的门禁 3 失败 —— 届时要么拆表要么改规格。
\providecommand{\tightlist}{}
\renewcommand{\tightlist}{\setlength{\itemsep}{2pt}\setlength{\parskip}{0pt}}
\AtBeginEnvironment{itemize}{%
  \setlength{\itemsep}{2pt}\setlength{\parsep}{0pt}%
  \setlength{\topsep}{4pt}\setlength{\partopsep}{0pt}}
\AtBeginEnvironment{enumerate}{%
  \setlength{\itemsep}{2pt}\setlength{\parsep}{0pt}%
  \setlength{\topsep}{4pt}\setlength{\partopsep}{0pt}}
\AtBeginEnvironment{description}{%
  \setlength{\itemsep}{2pt}\setlength{\parsep}{0pt}%
  \setlength{\topsep}{4pt}\setlength{\partopsep}{0pt}}
\setlength{\leftmargini}{1.5em}
\setlength{\leftmarginii}{1.2em}
\setlength{\leftmarginiii}{1.2em}

%% ---------------------------------------------------------------------
%% spec §11 硬约束的辅助：脚注（8/10）
%% ---------------------------------------------------------------------
\setlength{\footnotesep}{8pt}
\renewcommand{\footnoterule}{%
  \kern-3pt{\color{rule}\hrule width 0.4\columnwidth height 0.4pt}\kern2.6pt}

%% ---------------------------------------------------------------------
%% 链接：§3 accent 用于链接；PDF 书签层级 = 章 + 节
%% ---------------------------------------------------------------------
\AtBeginDocument{%
  \hypersetup{%
    colorlinks=true,%
    linkcolor=accent,%
    urlcolor=accent,%
    citecolor=accent,%
    filecolor=accent,%
    linktoc=page,%
    bookmarksnumbered=false,%
    bookmarksopen=true,%
    bookmarksopenlevel=0,%
    bookmarksdepth=1,%       章=0 节=1，§5 只收两级
    pdftitle={@BOOKTITLE@},%
    pdfborder={0 0 0},%
  }%
  %% §3 正文墨色（不用纯黑）。放在 AtBeginDocument 里，避免影响
  %% 宏包内部用色（caption 等已各自声明颜色）。
  \color{ink}%
  \Urlmuskip=0mu plus 1mu\relax
}

%% ---------------------------------------------------------------------
%% 标题页（规格未规定，最小实现：书名 22pt / 副题 14pt / 作者 / 日期）
%%   不抢版面：没有装饰线，靠留白。
%% ---------------------------------------------------------------------
\renewcommand{\maketitle}{%
  \begin{titlepage}%
    \thispagestyle{empty}%
    \null\vspace*{40mm}%
    \begin{flushleft}%
      {\fontsize{22}{28}\selectfont\heiti\bfseries\color{ink}@TITLE@\par}%
      \vspace{8pt}%
      {\fontsize{14}{18}\selectfont\color{muted}@SUBTITLE@\par}%
      \vspace{10mm}%
      {\color{accent}\rule{24mm}{0.8pt}\par}%
      \vspace{10mm}%
      {\fontsize{11}{15.2}\selectfont\color{muted}@AUTHOR@\par}%
      \vspace{2pt}%
      {\fontsize{9.5}{13}\selectfont\color{muted}@DATE@\par}%
    \end{flushleft}%
  \end{titlepage}%
}

%% ---------------------------------------------------------------------
%% 兼容层：pandoc 不同版本对 \pandocbounded / \tightlist / \subtitle 的处理
%% 不一致，这里统一兜底，避免 Undefined control sequence。
%% ---------------------------------------------------------------------
\providecommand{\subtitle}[1]{}
\providecommand{\pandocbounded}[1]{#1}
